\section{Reading the optimality conditions}\label{sec:sp:solution}

The system in Section~\ref{sec:sp:problem:foc} is short, but almost none of its content is on the
surface. This section works through what it says, in four steps: the pricing of materials
(Section~\ref{sec:sp:solution:materials}), what the accounting does \emph{not} price
(Section~\ref{sec:sp:solution:nowedge}), capital as a materials sink and the signs of the shadow
prices (Section~\ref{sec:sp:solution:signs}), and the recycling block, which contains the model's
central mechanism (Section~\ref{sec:sp:solution:recycling}). Section~\ref{sec:sp:solution:longrun}
then asks what the system does in the long run, and gets a more definite --- and more awkward ---
answer than one might expect.


\subsection{The pricing of materials}\label{sec:sp:solution:materials}

Everything in the materials block follows from one equation. Rearranging the extraction
condition~\eqref{eq:sp:foc:N},
\begin{align}
\boxed{\;B\;=\;p^N+d^Wp^P\;}
\qquad\text{that is}\qquad
\underbrace{F_R-p^W+d^Wp^P+\eta}_{\text{value of a secondary tonne}}
\;=\;
\underbrace{C^N_N+p^S+\Omega^{N,S}p^W}_{\text{cost of a virgin tonne}}
\;+\;d^Wp^P .
\label{eq:sp:solution:arbitrage}
\end{align}
Read from the left, this is a definition of $B$ that has been shown to equal something; read from
the right, it is the materials arbitrage condition. A tonne of material can be obtained in two ways.
Extracting it costs $C^N_N$ in resources, $p^S$ in foregone reserve value, and $\Omega^{N,S}p^W$ in
the waste cost of the overburden it drags up. Recovering it from waste instead avoids all three,
and additionally avoids the emission $d^W$ that the tonne would have caused had it stayed
untreated --- which is why $d^Wp^P$ appears on both sides and cancels out of the comparison of
\emph{sources} while remaining in the valuation of the recovered tonne.

The cancellation is the sharpest check available on the accounting. Note what has dropped out:
$p^W$, the shadow price of waste, appears on both sides of the arbitrage and is not what determines
whether recycling is worthwhile. It is worth being precise about why. A tonne of secondary material
does not \emph{remove} a tonne of waste from the economy; it postpones it. The tonne is recovered,
used in production, and returns as waste on the next pass. What recycling saves is not the waste
cost but the \emph{extraction} cost --- the resources, the rent, and the displaced overburden. This
is the substantive content of treating the materials balance as a ledger: a reduced-form model in
which recycling is credited with avoiding waste double-counts, because the waste is not avoided,
only deferred by one circuit.

Substituting~\eqref{eq:sp:solution:arbitrage} back into the definition of $B$ gives the companion
statement for the production margin,
\begin{align}
F_R &= p^N+p^W-\eta .
\label{eq:sp:solution:FR}
\end{align}
The marginal product of materials in production covers the full cost of delivering a virgin tonne
\emph{plus} the waste that tonne will eventually become. Here $p^W$ does appear, and it must: a
tonne drawn into the economy leaves as waste exactly once, whatever route it takes on the way, and
that is the point at which the waste cost is charged.


\subsection{What the accounting does not price}\label{sec:sp:solution:nowedge}

Two absences in Section~\ref{sec:sp:problem:foc} carry as much information as any of the equations.

\begin{proposition}[No wedge on output or consumption]\label{prop:sp:nowedge}
Under Version B the consumption condition is $u'(C)=\lambda$, with $\lambda$ the multiplier on the
goods constraint, and the marginal social value of a unit of output is $\lambda$. Neither carries a
correction for the waste that consumption or production generates.
\end{proposition}

This is not the statement that consumption and output generate no waste --- by
\eqref{eq:sp:setup:wc} they generate $\Omega^CC$ and $\Omega^YY$ tonnes of it. It is the statement
that at the margin the planner should not act on that fact. The reason is~\eqref{eq:sp:solution:FR}:
the mass that becomes $\Omega^CC$ was drawn into the economy as part of $R$ and was charged $p^W$ at
that point. Charging it again when it is spent on consumption would price the same tonne twice.

The contrast with a reduced-form specification is worth making explicit, because it is the main
methodological payoff of Section~\ref{sec:sp:setup:materialaccounting}. Suppose the ledger
\eqref{eq:sp:problem:ledger} were replaced by a posited waste function with \emph{fixed} intensities,
\begin{align}
W &= \Omega\left(\omega^YY+\omega^CC+\omega^{N}C^N+\omega^{D}C^D+\omega^WC^W\right)+\Phi^K\delta K,
&&\Omega \text{ a parameter},
\label{eq:sp:solution:reducedform}
\end{align}
which is~\eqref{eq:sp:setup:totalwaste_generic} with the substitution
\eqref{eq:sp:setup:ybalance_generic2} \emph{not} carried out. Redoing
Section~\ref{sec:sp:problem:foc} with \eqref{eq:sp:solution:reducedform} in place of the ledger
gives, in place of~\eqref{eq:sp:foc:C},
\begin{align}
u'\!\left(C\right) &= \lambda\left(1+\Omega^Cp^W\right),
&
\Omega^C&\equiv\Omega\omega^C,
\label{eq:sp:solution:spurious-C}
\end{align}
and a parallel wedge $\left(1-\Omega^Yp^W\right)$ on the marginal value of output --- a consumption
tax and an output tax, each proportional to the assumed intensity. These wedges are artefacts. They
appear because \eqref{eq:sp:solution:reducedform} lets the planner reduce economy-wide waste by
consuming less \emph{without reducing the mass drawn in}, which conservation of mass does not
permit. Once the accounting is done process by process, $\Omega$ is endogenous, the substitution
that produces the ledger cancels every $\omega^j$, and the wedges vanish. The difference is visible
in the optimal policy, not merely in the derivation: it is the difference between a model that
recommends taxing consumption on environmental grounds and one that does not.

\begin{proposition}[Irrelevance of the relative intensities]\label{prop:sp:irrelevance}
Under Version B, no $\omega^j$ and no $\Omega$ appears in any first-order
condition~\eqref{eq:sp:foc:C}--\eqref{eq:sp:foc:cs}, in any costate
equation~\eqref{eq:sp:costate:K}--\eqref{eq:sp:costate:M}, or in the constraint set of
Problem~\ref{prob:sp:planner}. The optimal allocation is therefore invariant to the relative
intensities, and $\Omega$ and $\Phi^Y$ are recovered afterwards
from~\eqref{eq:sp:setup:Omega} as reported quantities.
\end{proposition}

The proof is inspection, and the mechanism is the cancellation described after
\eqref{eq:sp:setup:ledger}: conservation of mass pins total waste down exactly, because every gram
of the accounted flow passed through $R$ on its way in, and how that gram was distributed across
activities on the way through does not change the total. The absolute intensities $\Omega^{N,S}$ and
$\Omega^{D,S}$ are a different matter --- they multiply flows that never passed through $R$, they
appear in~\eqref{eq:sp:foc:N} and~\eqref{eq:sp:foc:D}, and they do affect the allocation.

Proposition~\ref{prop:sp:irrelevance} has an uncomfortable corollary for the empirical side of the
project. Material-flow accounts report intensities, and it is natural to want to discipline the
$\omega^j$ against them; but under Version B no observation of the \emph{allocation} --- no
transition date, no recycling rate, no extraction path --- carries information about them. They are
identified only from the reported mass flows themselves, and they inform nothing else. Version A
reverses this.


\subsection{Capital as a materials sink, and the signs of the prices}\label{sec:sp:solution:signs}

The investment condition~\eqref{eq:sp:foc:I} is the one place where the stock-flow structure of
embodied material shows up in a price. Setting $\eta=0$ (the generic case, see below),
\begin{align}
q &= 1-\Phi^I\left(p^W-p^M\right).
\label{eq:sp:solution:q}
\end{align}
Building a unit of capital costs one unit of the final good and takes $\Phi^I$ tonnes out of the
current waste stream, worth $\Phi^Ip^W$; those tonnes return as waste over the life of the asset,
at a present value of $\Phi^Ip^M$. Capital is a materials sink, but a temporary one, and $q$ prices
the difference. Since~\eqref{eq:sp:costate:M} makes $p^M$ a discounted average of future $p^W$, the
bracket is positive whenever $p^W$ is positive and not falling fast, so $q<1$: the social cost of
installed capital is \emph{below} its goods cost. Along a path on which $p^W$ is constant,
\eqref{eq:sp:costate:M} gives $p^M=\delta p^W/\left(r+\delta\right)$ and
\begin{align}
q &= 1-\Phi^Ip^W\frac{r}{r+\delta},
\label{eq:sp:solution:q-stationary}
\end{align}
so the wedge is largest for durable capital ($\delta$ small) and vanishes as $\delta\to\infty$, when
capital stores nothing. It also vanishes as $r\to0$, when postponement is worthless.

\smalltitle{The sign of $p^W$.} Nothing guarantees $p^W>0$. From~\eqref{eq:sp:foc:W-solved},
\begin{align}
p^W<0
\qquad\Longleftrightarrow\qquad
\varpi\alpha\left(C^N_N+p^S+d^Wp^P\right)
\;>\;
c^W\!\left(\varpi\right)+\left(1-\varpi+d^W\varpi\right)p^P,
\label{eq:sp:solution:pW-sign}
\end{align}
that is, whenever the extraction a marginal tonne of waste displaces through recycling is worth more
than collecting, treating and emitting it costs. Since $p^S$ rises as reserves deplete
(Section~\ref{sec:sp:solution:longrun}) while $c^W$ and $p^P$ need not, this is not a pathological
case but a natural late-transition one: waste becomes an asset. Two consequences follow, and they
are worth recording because they are a useful internal check on any numerical solution.

First, by~\eqref{eq:sp:costate:M} the sign of $p^M$ follows the sign of $p^W$, so
by~\eqref{eq:sp:solution:q} the wedge on capital \emph{flips together with} the shadow price of
waste: when waste becomes an asset, embodying material in capital becomes a cost rather than a
saving, and $q>1$. The two flips are simultaneous and share a single cause, so a solution in which
they do not coincide is wrong. Second, any numerical implementation must therefore leave $p^W$, and
every instrument built from it, unbounded in sign.

\smalltitle{The multiplier $\eta$.} By~\eqref{eq:sp:foc:cs}, $\eta>0$ only where $\Phi^IG=R$: every
tonne entering production is embodied in new capital, and $\Omega=0$. This is a corner of the
accounting, not of the technology, and it is generic only in an economy investing at close to its
entire material throughput. Where it binds, $\eta$ enters $B$ and hence raises the value of both
virgin and secondary material by the same amount, as it must --- the constraint is on mass, and mass
is mass whatever its source. Under Version A the constraint disappears
(Section~\ref{sec:sp:versions:A}).

\smalltitle{The remaining signs.} Collecting~\eqref{eq:sp:costate:S}--\eqref{eq:sp:costate:M}: $p^S>0$
(reserves are an asset), $p^X<0$ (cumulative discovery is a liability), $p^P>0$ (pollution is a
liability, so the costate $-\lambda p^P$ is negative as posited), $p^M$ has the sign of $p^W$, and
$q$ has the opposite sign of the wedge. Only $p^W$ and $p^M$ are of ambiguous sign, and they are of
the same one.


\subsection{The recycling block}\label{sec:sp:solution:recycling}

The two conditions~\eqref{eq:sp:foc:varpi} and~\eqref{eq:sp:foc:KR} determine how circular the
economy is, and they respond to a single sufficient statistic, $B$.

\smalltitle{Recycling capital.} \eqref{eq:sp:foc:KR} reads $F_K=a'(x)B$: capital is allocated so
that its marginal product in production equals its marginal product in recycling, the latter being
$a'(x)$ tonnes of secondary material each worth $B$. This is a standard equalisation, but the
right-hand side is not standard: $B$ is a shadow price that depends on the scarcity of the resource,
so the recycling sector's demand for capital inherits the entire dynamics of the reserve.

\smalltitle{The treatment margin.} \eqref{eq:sp:foc:varpi} reads
$\alpha B+\left(1-d^W\right)p^P=dc^T/d\varpi$: treatment is pushed to where its marginal cost equals
the value of the material it yields plus the emission it prevents. The two terms are the two
distinct reasons to treat waste, and they are separately identified --- an economy with $d^W=1$
(treatment that does not immobilise anything) still treats waste, purely to recycle it; an economy
with $\alpha=0$ still treats waste, purely to avoid emissions.

\smalltitle{Circularity rises with scarcity.} The comparative statics of the block are unambiguous,
which is not obvious given that $\alpha$ and $a'$ move in opposite directions.

\begin{proposition}[Scarcity raises circularity]\label{prop:sp:circularity}
Hold $F_K$ and $p^P$ fixed and let $B$ rise. Then $x$ rises, $a(x)$ rises, the marginal yield
$\alpha$ rises, $\varpi$ rises, and the recirculated share $a\varpi$ rises.
\end{proposition}

\begin{proof}
From~\eqref{eq:sp:foc:KR}, $a'(x)=F_K/B$ falls, and $a''<0$ gives $dx/dB>0$; $a'>0$ then gives
$da/dB>0$. For the marginal yield, $\alpha=a(x)-a'(x)x$ has
$d\alpha/dx=a'-a''x-a'=-a''x>0$, so $\alpha$ rises with $x$ and hence with $B$. Therefore $\alpha B$
rises, so the left-hand side of~\eqref{eq:sp:foc:varpi} rises; $d^2c^T/d\varpi^2>0$
by~\eqref{eq:sp:setup:waste-cost} then gives $d\varpi/dB>0$. Both factors of $a\varpi$ rise.
\end{proof}

The mechanism of the paper is this proposition combined
with~\eqref{eq:sp:solution:arbitrage} and~\eqref{eq:sp:costate:S}. Depletion raises the in-situ rent
$p^S$; \eqref{eq:sp:solution:arbitrage} passes the rise one-for-one into $B$; and
Proposition~\ref{prop:sp:circularity} converts it into a higher recirculated share. By
\eqref{eq:sp:setup:ledger-solved} the circularity multiplier $\left(1-a\varpi\right)^{-1}$ rises with
it, so a given material throughput $R$ is supported by ever less virgin extraction. The transition
to a circular economy is, in this model, a scarcity-driven reallocation of capital, not a
pollution-driven one --- the pollution channel enters $B$ only through $d^Wp^P$, which is bounded by
the damage function, whereas $p^S$ is not bounded at all.

Note also what does \emph{not} appear in the transition mechanism: no $\omega^j$, no $\Omega$, no
$\Phi^Y$, and, by~\eqref{eq:sp:solution:arbitrage}, no $p^W$ either. The date at which the economy
turns circular is a function of extraction costs, resource rents, damages and the recycling
technology only. This is a strong prediction and a useful one, since it means the transition is not
sensitive to the least well identified part of the accounting.


\subsection{Stationarity and the long run}\label{sec:sp:solution:longrun}

It is natural to look for a steady state. The model as specified does not have one, and the reason
is instructive.

\begin{proposition}[No interior stationary point]\label{prop:sp:nosteadystate}
At any point where all five states of Problem~\ref{prob:sp:planner} are stationary, $D=N=0$ and
$W\left(1-a\varpi\right)=0$. Since $a(x)<1$ for all finite $x$, it follows that $W=R=R^R=0$: the
economy uses no materials.
\end{proposition}

\begin{proof}
$\dot X=D=0$ forces $D=0$, and then $\dot S=D-N=0$ forces $N=0$, and hence $\mathcal N=0$.
$\dot K=0$ gives $I=0$ and $G=\delta K$; $\dot M^K=0$ then gives $\Phi^IG=\delta M^K$, so the two
capital terms cancel in the ledger~\eqref{eq:sp:problem:ledger}, leaving $W=R$. But
$R=N+a\varpi W=a\varpi W$, so $W\left(1-a\varpi\right)=0$. With $\varpi\le1$ and $a<1$ the factor is
strictly positive.
\end{proof}

The culprit can be identified exactly. A stationary reserve with positive extraction requires
$D=N>0$, since $\dot S=D-N$; extraction must be replenished by discovery. But $\dot X=D$ makes
cumulative discovery a monotone state, so $D>0$ is itself incompatible with stationarity. The only
route to an interior rest point is therefore closed off by the single assumption that exploration
costs depend on cumulative past discovery. Drop it --- set $C^D_X=0$, so that $X$ ceases to be a
state --- and a genuine steady state with $D=N>0$ exists; this is switch~\ref{sw:sp:CDX} in
Section~\ref{sec:sp:versions:simple}, and it is the cheapest available repair. Note that removing
exploration altogether does \emph{not} repair it: with $\dot S=-N$, stationarity forces $N=0$
directly. The setup already flagged the corresponding structural feature --- there is no upper bound
on $S$, so the resource is never exhausted, only ever more costly --- and
Proposition~\ref{prop:sp:nosteadystate} is its dynamic counterpart: with $C^D_X>0$ the economy has no
resting point, it merely becomes asymptotically circular.

That is in fact the substantive content. Read the proposition in the limit rather than at a point.
As $t\to\infty$ the rent $p^S$ grows, $B$ grows with it by~\eqref{eq:sp:solution:arbitrage}, and
Proposition~\ref{prop:sp:circularity} sends $x\to\infty$, $a\to1$ and $\varpi\to1$. Along such a
path $N\to0$ while $R$ need not: from~\eqref{eq:sp:setup:ledger-solved} with $\dot M^K\to0$ and
$\mathcal N\to0$,
\begin{align}
R &= \frac{N-a\varpi\dot M^K+a\varpi\mathcal N}{1-a\varpi}\;\longrightarrow\;\frac{0}{0},
\label{eq:sp:solution:limit}
\end{align}
an indeterminate form whose value is whatever material throughput the recycling sector can sustain.
The emission flow~\eqref{eq:sp:problem:Xi} tends to $\left(1-\varpi+d^W\varpi\right)W-d^Wa\varpi W\to0$,
so the limiting economy is materially closed and emission-free, and it operates at positive scale.
This is the sense in which a fully circular economy is the model's long run: not a steady state that
is reached, but a boundary that is approached, and approached only because extraction becomes
infinitely expensive.

\smalltitle{The quasi-stationary block.} Even though the system has no rest point, the price block is
worth recording at $\dot{}=0$, both because it is what a slowly-moving path looks like locally and
because it is the natural initialisation for the numerical solution. Setting all price derivatives
to zero in~\eqref{eq:sp:costate:K}--\eqref{eq:sp:costate:M} and using $r=\rho$ from
$\dot\lambda=0$:
\begin{align}
F_K &= q\left(\rho+\delta\right),
&
p^S &= -\frac{C^N_S}{\rho},
&
p^X &= -\frac{C^D_X}{\rho},
\notag\\
p^P &= \frac{d}{\rho+\theta\!\left(P\right)+\theta'\!\left(P\right)P},
&
p^M &= \frac{\delta p^W}{\rho+\delta},
&
q &= 1-\Phi^Ip^W\frac{\rho}{\rho+\delta},
\label{eq:sp:solution:quasistationary}
\end{align}
with $p^W$ from~\eqref{eq:sp:foc:W-solved}, and the materially stationary ledger
\begin{align}
N+\mathcal N &= W\left(1-a\varpi\right)
\qquad\text{whenever }\dot M^K=0,
\label{eq:sp:solution:material-stationary}
\end{align}
which is~\eqref{eq:sp:setup:ledger-solved} at $\dot M^K=0$, and which requires $\Phi^K=\Phi^I$ by
\eqref{eq:sp:setup:phiK-motion}. Virgin extraction plus displaced mass equals waste not recirculated;
$N\to0$ as $a\varpi\to1$.

\smalltitle{A caveat on the multiplier.} \eqref{eq:sp:solution:material-stationary} makes the
quantitative stakes of the circularity multiplier visible. As $a\varpi\to1$ a given $N$ supports an
unbounded $W$, so the model's mass flows become arbitrarily sensitive to the recirculation
assumption in exactly the region the paper is about. The assumption in question is that recycled
material re-enters production contemporaneously --- $R^R=a\varpi W$ at the same instant --- which in
continuous time is innocuous but in a discrete-time implementation ties the number of circuits per
period to the period length. Section~\ref{sec:sp:versions} does not resolve this; it is flagged here
because $R/N$ and $W/N$ should be reported alongside any headline result, and their sensitivity to
the period length checked.


\subsection{The complete system}\label{sec:sp:solution:system}

For reference, and as the specification the quantitative implementation must reproduce, the
optimality system of Version B consists of:
\begin{itemize}[leftmargin=1.6em,itemsep=0.15em]
\item five state equations~\eqref{eq:sp:problem:states-motion} in $K,S,X,P,M^K$;
\item six price equations: \eqref{eq:sp:costate:K} for $q$, \eqref{eq:sp:costate:S} for $p^S$,
\eqref{eq:sp:costate:X} for $p^X$, \eqref{eq:sp:costate:P} for $p^P$, \eqref{eq:sp:costate:M} for
$p^M$, and \eqref{eq:sp:costate:euler} for $C$ (equivalently for $\lambda$);
\item five static conditions determining $I,N,D,\varpi,K^R$:
\eqref{eq:sp:foc:I}, \eqref{eq:sp:foc:N}, \eqref{eq:sp:foc:D}, \eqref{eq:sp:foc:varpi},
\eqref{eq:sp:foc:KR};
\item two definitions of prices: \eqref{eq:sp:foc:W} for $p^W$ and~\eqref{eq:sp:problem:damage} for
$d$;
\item the goods constraint~\eqref{eq:sp:problem:goods}, the ledger~\eqref{eq:sp:problem:ledger}, and
the complementary-slackness block~\eqref{eq:sp:foc:cs};
\item the definitions~\eqref{eq:sp:problem:definitions} for $R^R,R,Y,G$ and~\eqref{eq:sp:problem:Xi}
for $\Xi$;
\item initial conditions on the five states and the transversality
conditions~\eqref{eq:sp:problem:tvc}.
\end{itemize}
The accounting objects $\Omega$, $\Phi^Y$ and $\Phi^K$ are computed \emph{after} the fact,
from~\eqref{eq:sp:setup:Omega} and~\eqref{eq:sp:setup:phiK-motion}; by
Proposition~\ref{prop:sp:irrelevance} they do not feed back. This is the structural difference from
Version A, where they do, and where $\Omega$ has to be solved jointly with everything else.
